\section{A single public pair for the five component lower inequalities}
\label{app:same-pair}

\begin{proposition}[Level-2 same-public-pair lower]
\label{prop:same-pair-lower}
The displayed quantifiers are
\[
\forall \rho\text{ in the strict-interior field }0<v_-<B_0^2,\quad
\forall n\ge1,\quad
\forall(a,s,b,t)\in[0,1]^4,\quad
\exists(\Gamma^\star,P_X^\star),
\]
where \(\Gamma^\star=\Gamma^\star_{\rho,n,a,s,b,t}\) is one finite sequential-M1 public pair and \(P_X^\star\) is one covariate law, such that
\[
\inf_{\delta_{\Gamma^\star}\ {\rm Borel}}
\sup_{P\in\cP(\Gamma^\star)}
\E_P\abs{\delta_{\Gamma^\star}(Z_{1:2n})-\beta(P)}
\ge c_\rho^\star(q+X)
\]
for a constant \(c_\rho^\star>0\) depending only on fixed \(\rho\).  Its compatible-law family contains \(q\) and every active native P/O/\(ab\)/M family; inactive component inequalities use the common \(q\) splice.  Thus the five inequalities hold on one covariate support and marginal \(P_X^\star\), without an explicit component tag.  The pair is fixed before Nature chooses component, coordinate, sign, label, or data, and the inequality covers every arbitrary Borel rule that knows the entire pair.
\end{proposition}

This appendix proves Proposition~\ref{prop:same-pair-lower}.  It reuses the
native nonparametric raw laws proved in Appendices~\ref{app:qp}--\ref{app:m}
and constructs the following common-host q variant, but not an
external certificate: we define the enlarged pair, verify its complete
sequential-M1 and budget interface, check that the component laws remain in
the same compatible-law set, and apply testing directly to a rule that knows
the enlarged pair.

\subsection{Host, active branches, and public pair}

Fix \(0<v_-<B_0^2\), \(n\ge1\), and
\((a,s,b,t)\in[0,1]^4\).  Write \(N=2n\), \(v_0=v_-\), and
\[
 m_\rho=\min\left\{\frac1{64},\frac{v_0}{4B_0},
                \frac{B_0-\sqrt{v_0}}4\right\},\qquad
 \lambda_\rho=\min\left\{1,\frac{192c_{\rm rad}}{23},
                128B_{\mathcal G}\right\}.
\]
Both constants are strictly positive and depend only on fixed \(\rho\).
Declare a large P, O, \(ab\), or M branch, respectively, when
\[
\begin{array}{c|c|c}
\text{branch}&\text{large-branch predicate}&\text{dictionary size}\\ \hline
P&n\ge16,\ s^2\ge q,\ s(b\wedge s)>0&
 p_P=\lceil\exp(64Ns^2)\rceil\\
O&n\ge16,\ t^2\ge q,\ t(a\wedge t)>0&
 p_O=\lceil\exp(64Nt^2)\rceil\\
ab&n\ge16,\ a^2+b^2\ge q,\ ab>0&
 p_{ab}=\lceil\exp(64N(a^2+b^2))\rceil\\
M&n\ge16,\ (s\wedge t)^2\ge q,\ s\wedge t>0&
 p_M=\lceil\exp(64N(s\wedge t)^2)\rceil .
\end{array}
\]
An inactive branch has count zero.  Put
\[
 K=p_P+p_O+p_{ab}+p_M,\qquad
 m_\star=\max\{1,\lceil\log_2(K+1)\rceil\}.
\]
The group \((\mathbb Z/2\mathbb Z)^{m_\star}\) has at least \(K\)
distinct nonzero Walsh characters.  Partition \(K\) of them into four
disjoint frequency sets
\(\mathcal J_P,\mathcal J_O,\mathcal J_{ab},\mathcal J_M\) of the
displayed sizes.  By the fixed-infinite-space convention of
Appendix~\ref{app:model}, preselect a measurable injection of this finite
group into \(\mathcal X\), extend every character by zero outside the
image, and let \(P_X^\star\) be uniform on the whole image.  Thus every
allocated character \(\phi_j\) is Rademacher-valued, mean zero, and
unit-norm under \(P_X^\star\), and distinct frequencies are orthogonal.
The support and \(P_X^\star\) are identical for every component, coordinate,
sign, and label.  No component coordinate or disjoint-support event is
observed.

For an active P branch, let \(w_P=b\wedge s\) and set
\[
 \eta_P=\frac{m_\rho w_P}{4},\quad
 \gamma_P=\frac{\lambda_\rho s}{256},\quad
 \Delta_P=\frac{2\eta_P\gamma_P}{v_0},\quad
 \kappa_P=\gamma_P+\Delta_P\eta_P.
\]
For active O and M branches, put
\[
\begin{aligned}
 w_O&=a\wedge t,&
 \gamma_O&=\frac{\lambda_\rho m_\rho t}{4},&
 \eta_O&=\frac{w_O}{256},&
 \Delta_O&=\frac{\gamma_O\eta_O}{v_0+\gamma_O^2},\\
 w_M&=s\wedge t,&
 \gamma_M&=\frac{\lambda_\rho m_\rho w_M}{4},&
 \eta_M&=\frac{\lambda_\rho w_M}{256},&
 \Delta_M&=\frac{\gamma_M\eta_M}{v_0+\gamma_M^2}.
\end{aligned}
\]
For an active \(ab\) branch, retain
\[
 \eta_{ab}=\frac{m_\rho b}{4},\quad
 \gamma_{ab}=\frac{a}{512},\quad
 \Delta_{ab}=\frac{2\eta_{ab}\gamma_{ab}}{v_0},\quad
 \kappa_{ab}=\gamma_{ab}+\Delta_{ab}\eta_{ab}.
\]
Terms from inactive branches are omitted from the following definitions.

% C6-OUTCOME-CLASS-BEGIN
\[
\mathcal G_\mu^\star
=\{0\}
\cup\{\sigma\gamma_P\phi_j,\sigma\kappa_P\phi_j:
          j\in\mathcal J_P,\ \sigma\in\{-1,1\}\}
\cup\{\sigma\eta_M\phi_j:
          j\in\mathcal J_M,\ \sigma\in\{-1,1\}\}.
\]
% C6-OUTCOME-CLASS-END
% C6-PROPENSITY-CLASS-BEGIN
\[
\mathcal G_\pi^\star
=\{0\}
\cup\{\sigma\gamma_O\phi_j:
          j\in\mathcal J_O,\ \sigma\in\{-1,1\}\}
\cup\{\sigma\gamma_M\phi_j:
          j\in\mathcal J_M,\ \sigma\in\{-1,1\}\}.
\]
% C6-PROPENSITY-CLASS-END
Fix ordered countable representations
\(\mathbf g_\mu^\star=(g_{\mu,j}^\star)_{j\in\mathbb N}\) and
\(\mathbf g_\pi^\star=(g_{\pi,j}^\star)_{j\in\mathbb N}\) by listing every
element of the corresponding finite class and then repeating zero, and set
\[
 \Gamma^\star=(\mathcal G_\mu^\star,\mathcal G_\pi^\star,
                 \mathbf g_\mu^\star,\mathbf g_\pi^\star).
\]
Each class is a finite union, not a linear span.  Taking spans would permit
additive combinations across bands, change the pointwise envelope, and make
the following finite-cardinality complexity certificate inapplicable.
The \(ab\) band is deliberately absent from both public classes because its
size is controlled by \(a^2+b^2\), not by \(s^2\) or \(t^2\).

\subsection{Sequential-M1, envelopes, and stochastic budgets}

These deterministic enumerations are sequential M1 representations; their
measurable evaluation maps are well-defined because the embedded finite-block
characters are measurable.  Pointwise convergence and the common candidate
envelope are immediate.  The component bounds in Appendices~\ref{app:qp} and
\ref{app:m} give
\[
 \kappa_P\le2\gamma_P\le\frac{\lambda_\rho s}{128}\le B_{\mathcal G},
 \qquad
 \eta_M,\gamma_O,\gamma_M\le\frac{\lambda_\rho}{256}\le B_{\mathcal G}.
\]
Because the classes are unions rather than sums, their envelopes are the
largest displayed amplitudes.

Suppose the outcome class is nontrivial and put
\(p_s=\lceil\exp(64Ns^2)\rceil\).  Active M implies \(p_M\le p_s\), while
active P gives \(p_P=p_s\).  Counting the full difference class, including
cross-band differences, yields
\[
 |\mathcal G_\mu^\star|\le1+4p_P+2p_M\le7p_s,\qquad
 |\Delta\mathcal G_\mu^\star|\le49p_s^2
 \le196\exp(128Ns^2).
\]
Nontriviality implies \(Ns^2\ge8\); hence
\(\log|\Delta\mathcal G_\mu^\star|<129Ns^2=258ns^2\).
The difference-class envelope is at most twice
\(H_\mu=\max\{\kappa_P,\eta_M\}\le\lambda_\rho s/128\).
The finite-class Massart inequality from Lemma~\ref{lem:massart}, applied
globally and hence to every localized subset, gives
\[
 \mathfrak R_n(r;\Delta\mathcal G_\mu^\star)
 <46H_\mu s\le\frac{23\lambda_\rho}{64}s^2,
 \qquad
 \frac{\mathfrak R_n(r;\Delta\mathcal G_\mu^\star)}{3r}
 \le\frac{23\lambda_\rho s}{192}\le c_{\rm rad}s
\]
for every \(r\ge s>0\).

Similarly, with \(p_t=\lceil\exp(64Nt^2)\rceil\),
\[
 |\mathcal G_\pi^\star|\le1+2p_O+2p_M\le5p_t,\qquad
 |\Delta\mathcal G_\pi^\star|\le25p_t^2
 \le100\exp(128Nt^2).
\]
Nontriviality implies
\(\log|\Delta\mathcal G_\pi^\star|<258nt^2\), and the difference envelope
is twice \(H_\pi\le\lambda_\rho t/256\).  Therefore
\[
 \mathfrak R_n(r;\Delta\mathcal G_\pi^\star)
 <46H_\pi t\le\frac{23\lambda_\rho}{128}t^2,
 \qquad
 \frac{\mathfrak R_n(r;\Delta\mathcal G_\pi^\star)}{3r}
 \le\frac{23\lambda_\rho t}{384}\le c_{\rm rad}t
\]
for every \(r\ge t>0\).  If \(s=0\), neither P nor M is active, so
\(\mathcal G_\mu^\star=\{0\}\) when \(s=0\); if \(t=0\), neither O nor
M is active, so \(\mathcal G_\pi^\star=\{0\}\) when \(t=0\).  The
corresponding difference class is then exactly zero at every positive
radius, and no radius-zero quotient is used.

\subsection{Approximation and raw-law membership}

For an off-class truth \(c\phi_j\) and a nonzero public candidate
\(d\phi_k\) from another band,
\[
 \|c\phi_j-d\phi_k\|_2^2=c^2+d^2>c^2
 =\|c\phi_j\|_2^2.
\]
Consequently the added bands preserve the P/O/ab approximation distances
exactly.  In detail, q has two zero truths; P has its outcome truth in
\(\mathcal G_\mu^\star\) and propensity distance \(\eta_P\le b\);
O has its propensity truth in \(\mathcal G_\pi^\star\) and outcome distance
\(\eta_O\le a\); \(ab\) uses \(\mathcal J_{ab}\) off both public classes
and has the original distances bounded by \(a,b\); M has both truths in
the corresponding public classes.

We now audit the q/P/O/ab/M raw laws.  For q, use \(X\sim P_X^\star\)
independently of
\[
 T\in\{-L,0,L\},\qquad L=\sqrt{v_0}+2m_\rho,\qquad
 \E[T]=0,\quad \E[T^2]=v_0,
\]
with explicit masses
\[
 \Pr(T=L)=\Pr(T=-L)=\frac{v_0}{2L^2},\qquad
 \Pr(T=0)=1-\frac{v_0}{L^2}.
\]
Because \(L^2\ge v_0\), these probabilities are nonnegative.  They sum to
one since \(v_0/L^2+1-v_0/L^2=1\), and
\[
 \E T=L\,v_0/(2L^2)-L\,v_0/(2L^2)=0,
 \qquad
 \E T^2=L^2v_0/L^2=v_0.
\]
Thus \(\E T=0\) and \(\E T^2=v_0\) exactly.  Set the target
\(\theta=\pm d_\rho^\star/\sqrt n\), where
\[
 d_\rho^\star=\min\left\{3,\frac{B_0}{16},
                  \frac{B_0}{16\sqrt{v_0}}\right\}.
\]
Conditional on \((X,T)=(x,t)\), let \(R\in\{-1,+1\}\) satisfy
\[
 \Pr(R=r\mid X,T=t)=\frac12\left(1+\frac{2r\theta t}{B_0}\right),
 \qquad r\in\{-1,+1\},
\]
and set \(Y=B_0R/2\).
This is the three-point common-host variant of the scalar q experiment in
Appendix~\ref{app:q-lower}.  The bounds defining \(m_\rho\) imply
\[
 L\le\frac{B_0+\sqrt{v_0}}2<B_0,\qquad
 \frac{L}{\sqrt{v_0}}\le1+\frac{\sqrt{v_0}}{2B_0}<\frac32,
 \qquad \abs{\frac{2\theta T}{B_0}}\le\frac3{16}.
\]
Thus the two response probabilities lie in \([13/32,19/32]\),
\(\E(R\mid X,T)=2\theta T/B_0\), and
\(\E(Y\mid X,T)=\theta T\).  With zero nuisances,
\(\varepsilon=Y-\theta T\) is conditionally centered and obeys
\(\abs\varepsilon\le B_0/2+3B_0/32<B_0\); the target, treatment, outcome,
and treatment residual obey the same shared envelope, and the latter has
constant variance \(v_0\).  With this displayed response likelihood held
fixed, the likelihood-energy calculation's dependence on the treatment
marginal uses only this second moment \(\E T^2=v_0\).  The response likelihood
itself is still used.  Consequently its \(K_q\), testing constant, and rate
are unchanged.

For P, O, \(ab\), and M, use respectively the raw conditional masses from
Appendices~\ref{app:p-raw}, \ref{app:o}, \ref{app:ab}, and \ref{app:m},
replacing only their Walsh coordinate by a coordinate in
\(\mathcal J_P,\mathcal J_O,\mathcal J_{ab},\mathcal J_M\).  A nonzero
character on the larger host is still a mean-zero unit Rademacher variable,
so every atom probability, normalization identity, conditional-centering
identity, response-centering identity, target range, treatment/response/
nuisance/residual envelope, and constant conditional variance calculation
is unchanged.  In particular the common reference laws remain valid because
\[
 \Delta(v_0+\eta^2)=\eta(\kappa+\gamma)
 \quad\text{for P and \(ab\)},\qquad
 \Delta(v_0+\gamma^2)=\gamma\eta
 \quad\text{for O and M}.
\]
Thus the \(q\) family and every active native P/O/\(ab\)/M hard subfamily lie
in the single \(\mathcal P(\Gamma^\star)\).  An inactive component inequality
reuses the common \(q\) splice; hence all five component lower inequalities
hold within the pair even though the named subfamilies need not be five
distinct laws at every cell.  The proof uses strict lower-endpoint slack
only through \(m_\rho>0\) and the component atom envelopes; constant
conditional variance is inherited from the raw laws; the fixed infinite
ambient space is used only to embed the finite host; and known pre-data upper certificates are
used only to choose the bands and amplitudes before data.

\subsection{Likelihood distance and arbitrary-rule reduction}

For two different coordinates in any component band,
\(\E_{P_X^\star}\phi_j\phi_k=0\); equal coordinates with equal or opposite
signs give the same \(1\pm K_\ell\) likelihood inner products as in the
component appendices.  Hence their exact finite Gram-mixture
\(\chi^2\)/Hellinger bounds, and therefore their product-law total-variation
bounds, are unchanged.  For q, the exact one-sample likelihood energy is
\[
 K_q=\frac{4(d_\rho^\star)^2v_0}{nB_0^2}\le\frac1{64n},
\]
so the original \(2n\)-sample testing threshold also remains strict.

Fix the fully disclosed \(\Gamma^\star\), frequency allocation, \(P_X^\star\),
and budgets.  Any allowed \(\delta_{\Gamma^\star}\) is now an ordinary Borel
statistic of \(Z_{1:2n}\).  The two-mixture expected-loss inequality applies
to every arbitrary Borel rule after disclosure of the full pair:
\[
 \max_{\ell\in\{0,1\}}
 \E_{\Pi_{i,\ell}}\!
 \left|\delta_{\Gamma^\star}-\theta_{i,\ell}\right|
 \ge \frac{\Delta_i}{4}
 \{1-\operatorname{TV}(\Pi_{i,0},\Pi_{i,1})\}.
\]
Every \(\Pi_{i,\ell}\) is supported on a subfamily of the same
\(\mathcal P(\Gamma^\star)\).  Evaluating unused public functions at observed
\(X\) is deterministic post-processing of \(X\); it is already available to
the arbitrary statistic and creates no label-dependent side channel.  Thus
the richer dictionary cannot invalidate a certified testing distance.

\subsection{All-sample-size splice and fixed-pair assembly}

When a large branch is inactive, its named splice subfamily is the same q
subfamily inside the same pair.  If \(s^2<q\), then
\(P=s(b\wedge s)\le s^2<q\); if \(n<16\), then \(P\le1\le4q\).
The O statement is the exact mirror.  If \(a^2+b^2<q\), then
\(ab\le(a^2+b^2)/2<q/2\); if \(n<16\), then \(ab\le1\le4q\).
If \((s\wedge t)^2<q\), then \(M<q\); if \(n<16\), then
\(M\le1\le4q\).  The equality boundaries belong to the large branch by
definition, and zero terms require no division or limiting argument.

Use \(c_{\min,\rho}\) from Appendix~\ref{app:assembly}, the minimum of the q
testing constant, the four fixed large-branch constants, and all four q-splice
constants.  It depends only on fixed
\(\rho\), not on \(n,a,s,b,t,\Gamma^\star\), or a compatible law.  Because
all five inequalities hold for every rule on the same pair, the infimum is
taken only after their maximum:
\[
\begin{split}
 \inf_{\delta_{\Gamma^\star}\ {\rm Borel}}
 \sup_{P\in\mathcal P(\Gamma^\star)}
 \E_P|\delta_{\Gamma^\star}-\beta(P)|
 &\ge c_{\min,\rho}\max\{q,ab,M,O,P\}\\
 &\ge\frac{c_{\min,\rho}}5(q+S_0)
 \ge\frac{c_{\min,\rho}}{10}(q+X),
\end{split}
\]
where \(S_0=ab+M+O+P\) and \(S_0\le X\le2S_0\) by
\eqref{eq:budget-equivalence}.  Setting
\(c_\rho^\star=c_{\min,\rho}/10\) proves
Proposition~\ref{prop:same-pair-lower}.  This is maximum-to-sum within one fixed
public pair.  Taking the outer supremum in \eqref{eq:risk} then proves the
I-stratum lower of Theorem~\ref{thm:main}, but no change of pair is used in
the substantive argument.

\subsection{Every prescribed atomless covariate marginal}
\label{app:fixed-nu}

Let \(\nu\) be a probability measure on the fixed standard Borel space
\(\mathcal X\), prescribed before the public pair is chosen.  Define
\[
 \mathcal P_\nu(\Gamma)
 =\{P\in\mathcal P(\Gamma):P_X=\nu\},\qquad
 \mathfrak G_\nu=\{\Gamma\text{ legal}:\mathcal P_\nu(\Gamma)\ne\varnothing\},
\]
and define the fixed-marginal experiment by
\begin{equation}
 \mathfrak R_{\nu,n}(a,s,b,t)
 =\sup_{\Gamma\in\mathfrak G_\nu}
   \inf_{\delta_{\Gamma,\nu}\ {\rm Borel}}
   \sup_{P\in\mathcal P_\nu(\Gamma)}
   \E_P\left|\delta_{\Gamma,\nu}(Z_{1:2n})-\beta(P)\right|.
 \label{eq:risk-fixed-nu}
\end{equation}
The learner may know \(\nu\) as well as the entire represented pair.
Nonemptiness here is with respect to this prescribed marginal, rather than
compatibility with some other covariate law.

\begin{proposition}[Same-pair lower for every prescribed atomless marginal]
\label{prop:fixed-nu-lower}
For each fixed strict-interior \(\rho\), there are constants
\(0<c_\rho\le C_\rho<\infty\), independent of \(\nu\), such that
\[
 \forall\nu\text{ atomless on }\mathcal X,\quad
 \forall n\ge1,\quad\forall(a,s,b,t)\in[0,1]^4,\quad
 \exists\Gamma^\star_{\nu,\rho,n,a,s,b,t}\in\mathfrak G_\nu
\]
with finite sequential-M1 classes and
\begin{equation}
 \inf_{\delta_{\Gamma^\star,\nu}\ {\rm Borel}}
 \sup_{P\in\mathcal P_\nu(\Gamma^\star)}
 \E_P|\delta_{\Gamma^\star,\nu}-\beta(P)|
 \ge c_\rho(q+X).
 \label{eq:fixed-nu-same-pair}
\end{equation}
This single compatible-law set contains the q family and all active native
P/O/\(ab\)/M families, and supplies the inactive inequalities by the q
splice.  Consequently
\[
 c_\rho(q+X)\le\mathfrak R_{\nu,n}(a,s,b,t)
 \le C_\rho(q+X).
\]
At saturation \(v_-=B_0^2\), for every prescribed marginal \(\nu\),
including atomic marginals,
\[
 \frac{3}{128\sqrt n}\le\mathfrak R_{\nu,n}(a,s,b,t)
 \le\frac1{\sqrt{2n}}.
\]
If \(v_->B_0^2\), every \(\mathcal P_\nu(\Gamma)\) is empty.
\end{proposition}

\begin{proof}
We first construct measurable Walsh functions for the prescribed atomless
\(\nu\).  A standard Borel space has a countable Borel family
\((A_k)_{k\ge1}\) separating its points: take the pullbacks of a countable
base under a Borel realization in a Polish space.  The map
\[
 e(x)=\sum_{k=1}^{\infty}2\,3^{-k}\mathbf1_{A_k}(x)\in[0,1]
\]
is Borel as a limit of Borel partial sums and is injective.  Indeed, at the
first differing digit \(k\), the difference \(2\,3^{-k}\) exceeds the
largest possible opposing tail \(\sum_{j>k}2\,3^{-j}=3^{-k}\).
Thus \(\mu=e_\#\nu\) has no atoms: every fiber of \(e\) is empty or a
singleton.  Its distribution function
\(F(z)=\mu(( -\infty,z])\) is continuous, since the jump at \(z\) equals
\(\mu\{z\}=0\).  Set
\[
 U(x)=F(e(x)).
\]
This is Borel and \(U_\#\nu\) is uniform on \([0,1]\).  To see the
probability integral transform directly, for \(0<t<1\) continuity and the
limits of \(F\) give a point with \(F(z)=t\).  The right endpoint
\(b_t=\sup\{z:F(z)\le t\}\) is finite, satisfies \(F(b_t)=t\), and
\(\{z:F(z)\le t\}=(-\infty,b_t]\).  Hence
\(\nu\{U\le t\}=F(b_t)=t\); the endpoints follow by limits.

For \(0\le u<1\) define
\[
 b_k(u)=\lfloor2^ku\rfloor-2\lfloor2^{k-1}u\rfloor,\qquad
 r_k(u)=(-1)^{b_k(u)},
\]
and set \(b_k(1)=0\) and \(r_k(1)=1\).  These Borel bits are defined even at dyadic
endpoints; our convention changes no distributional identity.  The first
\(m\) bits have each binary pattern on one half-open interval of length
\(2^{-m}\), so their vector is uniform on \(\{0,1\}^m\).
Enumerate the nonempty finite subsets \(H\subset\mathbb N\) and put
\[
 w_H(u)=\prod_{k\in H}r_k(u),\qquad \phi_H(x)=w_H(U(x)).
\]
Every \(\phi_H\) is Borel, everywhere sign-valued, centered under \(\nu\),
and has squared norm one.  For distinct \(H,J\),
\(\phi_H\phi_J=\phi_{H\triangle J}\), so
\(\E_\nu\phi_H\phi_J=0\).  These characters are orthogonal; we do not
require them to be jointly independent.

Use precisely the active predicates, band counts, \(m_\star\), amplitudes,
and finite unions already defined in this appendix.  Allocate the bands
among the nonempty subsets of \(\{1,\ldots,m_\star\}\), identifying such
a subset with its nonzero binary frequency.  Replace each finite-host
character by the corresponding \(\phi_H\).  The entire finite vector of
allocated characters under \(\nu\) then has exactly the same joint law
as under the earlier uniform finite host, because the first \(m_\star\)
bits are uniform.  Define \(\Gamma^\star_{\nu,\rho,n,a,s,b,t}\) using the
same ordered finite lists followed by repetitions of zero.  This is a
public choice made before the component, coordinate, sign, and label.

All functions are measurable, and these representations satisfy sequential
M1 with their finite classes.  Their pointwise envelopes remain
\(H_\mu\le\lambda_\rho s/128\) and
\(H_\pi\le\lambda_\rho t/256\).  Orthogonality gives exactly the same
off-band approximation calculation
\(\|c\phi_H-d\phi_J\|_{L^2(\nu)}^2=c^2+d^2\) for \(H\ne J\).
Thus the P/O/\(ab\) distances are preserved, the M truths remain in
class, and the q truths remain zero.  The complete difference classes,
including cross-band differences, still obey
\[
 |\Delta\mathcal G_\mu^\star|\le49p_s^2,\qquad
 |\Delta\mathcal G_\pi^\star|\le25p_t^2
\]
whenever the corresponding class is nontrivial.  The same global Massart
calculation therefore gives, at every positive radius in the stated ranges,
\[
 \begin{aligned}
 \frac{\mathfrak R_n(r;\Delta\mathcal G_\mu^\star)}{3r}
 &\le\frac{23\lambda_\rho s}{192}\le c_{\rm rad}s
 &&(r\ge s>0),\\
 \frac{\mathfrak R_n(r;\Delta\mathcal G_\pi^\star)}{3r}
 &\le\frac{23\lambda_\rho t}{384}\le c_{\rm rad}t
 &&(r\ge t>0).
 \end{aligned}
\]
Zero stochastic budgets again give singleton zero classes, with zero
complexity at every positive radius.  There is no dependence on \(\nu\)
in these bounds.

For each active component, retain the original conditional kernel of
\((T,Y)\) given its signed Walsh value, and let \(X\sim\nu\).
This is a Borel probability kernel because it has finitely many atoms with
Borel masses.  At each \(x\), the signed value is either \(-1\) or \(1\),
so the original atom positivity, normalization, target range, all shared
envelopes, both conditional-centering identities, and exact constant
treatment-residual variance \(v_0\) hold pointwise.  In particular this
argument preserves the raw laws, rather than only their reduced moments.
Together with the class calculations it proves membership in
\(\mathcal P_\nu(\Gamma^\star)\).

We check testing on the full observation space.  For each active component
\(i\in\{P,O,ab,M\}\), let \(Q_i\) be its native finite-atom reference
law on \((T,Y)\).  A different component may use a different \(Q_i\), but
the two labels within component \(i\) use the same \(Q_i\), by the
reference identities verified above.  On
\(\mathcal X\times\mathbb R^2\) take
\(\widetilde Q_i=\nu\otimes Q_i\).  The full one-observation likelihood
is exactly
\[
 \frac{dP_{i,\ell,j,\sigma}}{d\widetilde Q_i}(x,t,y)
 =1+\sigma\phi_j(x)h_{i,\ell}(t,y),\qquad
 \E_{Q_i}h_{i,\ell}=0,\quad
 K_{i,\ell}=\E_{Q_i}h_{i,\ell}^2.
\]
Consequently, for a fixed label, its full Gram entries are
\[
 \E_{\widetilde Q_i}
 [(1+\sigma\phi_jh_{i,\ell})(1+\tau\phi_kh_{i,\ell})]
 =1+\sigma\tau K_{i,\ell}\E_\nu(\phi_j\phi_k).
\]
They equal \(1\) off the coordinate diagonal and \(1\pm K_{i,\ell}\)
on it.  For the uniform coordinate/sign mixture of \(N=2n\) independent
observations, Lemma~\ref{lem:testing-tools} gives exactly
\[
 \chi^2(\overline P_{i,\ell}^{\,N},\widetilde Q_i^{\,N})
 =\frac{\{(1+K_{i,\ell})^N+(1-K_{i,\ell})^N\}/2-1}{p_i}.
\]
The native energy bounds and all ensuing Hellinger and total-variation
bounds are therefore unchanged.  This calculation is for the full \(X\),
not a projection onto its first bits.  Extra digits and evaluations of all
public functions are measurable information already included in this
experiment.  Thus an arbitrary Borel learner knowing both \(\nu\) and the
full pair is covered by the same testing inequality.

For q use the three-point variant proved in this appendix, independently
of \(X\sim\nu\), with \(\theta=\pm d_\rho^\star/\sqrt n\).
The zero truths belong to the enlarged classes.  Its full reference is
\(\nu\otimes Q_q\), where \(Q_q\) has the displayed three-point
treatment marginal and an independent uniform response sign.  Its
likelihood is \(1\pm2d_\rho^\star rt/(B_0\sqrt n)\), with unchanged
energy \(K_q\le1/(64n)\).  Hence the same q lower constant holds.
This also proves \(\mathcal P_\nu(\Gamma^\star)\ne\varnothing\) for
every budget cell.  All inactive-branch comparisons with q are numerical
and unchanged, including \(n<16\), zero budgets, and equality boundaries.
Taking the maximum of the five within-pair inequalities before taking the
infimum over Borel rules gives
\[
 \inf_\delta\sup_{P\in\mathcal P_\nu(\Gamma^\star)}
 \E_P|\delta-\beta(P)|
 \ge\frac{c_{\min,\rho}}{10}(q+X).
\]
This proves \eqref{eq:fixed-nu-same-pair} with the same constant as for the
finite host.  The global upper of Proposition~\ref{prop:upper} applies to
each restricted compatible-law set, and its pair-specific Borel estimator
is allowed in \eqref{eq:risk-fixed-nu}; the upper constant is likewise
independent of \(\nu\).

For saturation, the pairwise upper (H.1)--(H.4) requires no assumption on
the marginal.  In the lower experiment of Appendix~\ref{app:boundary},
replace only its constant covariate by an independent \(X\sim\nu\).
Keep the singleton zero classes, \(T=\pm B_0\) equiprobably,
\(\delta_n=3/(32\sqrt n)\), and exactly the response probabilities
in (H.5).  All raw-law and zero-budget checks remain valid.  The common
factor \(\nu\) contributes affinity one, so the product-law calculation
still gives \(\operatorname{TV}^2\le72/1024<1/4\), and (H.6) gives
\(3/(128\sqrt n)\).  These saturation constants are not obtained by
continuing the interior constants.  Finally the envelope contradiction
when \(v_->B_0^2\) holds for every marginal.
\end{proof}

The interior strengthening is not claimed for arbitrary atomic marginals.
For example,
if \(\nu\) is a Dirac mass, the model reduces to bounded scalar regression
with an intercept, and \(\beta=\operatorname{Cov}(T,Y)/\operatorname{Var}(T)\)
with \(\operatorname{Var}(T)\ge v_-\).  Clipping the empirical covariance
divided by the empirical variance truncated below at \(v_-/2\) gives
uniform expected error \(O_\rho(q)\): bounded empirical first and second
moments have expected error \(O_\rho(q)\), the truncated denominator is
bounded away from zero, and truncation cannot increase its error from the
true variance.  The q family gives the matching lower.  Thus fixed
positive nuisance budgets cannot force the interior nonparametric terms
for this marginal.  The original finite-host proposition remains valid
on every fixed infinite standard Borel host, including countably infinite
hosts that admit no atomless probability measure.

\subsection{Scope of the strengthening}

The construction is Level-2 common-support multiplexing.  It is not an
observable tagged direct sum, because all laws have the same \(X\) marginal
on the same support and cross-band differences are paid in the public
complexity.  It is also not Level 3: \(\Gamma^\star\) is engineered, may
depend on public \(n\) and the budget cell, and supplies no unknown-budget,
common-rule, computational, or unrestricted heteroskedastic extension.
The original finite-host result applies to countably infinite hosts; the
fixed-marginal interior result assumes an atomless marginal.  Walsh multiplexing is a proof construction, not an algorithmic
design recommendation or a generic lower-bound claim.
